\documentclass{report}
\usepackage{amsmath,amssymb}
\setlength{\parindent}{0mm}
\setlength{\parskip}{1em}
\begin{document}
\begin{center}
\rule{6in}{1pt} \
{\large Fande, FD Kong \\
{\bf A scalable parallel nonlinear solver for numerical simulation of fluid-structure interactions with applications in hemodynamics}}

Department of Computer Science \\ University of Colorado Boulder \\ Boulder \\ Colorado 80309
\\
{\tt fande.kong@colorado.edu}\\
Xiao-Chuan/XC Cai\end{center}

Numerical simulation of fluid-structure interactions (FSI) with a
patient-specific geometry provides a valuable tool in computer-aided
surgeries. Especially, with the rapid development of
supercomputers, a high resolution of the numerical simulation of cardiovascular dynamics
becomes achievable, and calculated metrics such as the blood velocity, the
blood pressure and the wall displacement can be used directly or indirectly as valuable
means, instead of empirical and often risky clinical experiments, to
examine the suitability
and the efficiency of various reconstructive procedures. However, the
parallel numerical solution of the FSI system with a patient-specific
geometry is still challenging
because (1) computation domains are irregular and deformable so that the
construction of efficient solvers such as
multigrid methods or multilevel methods is difficult; (2) fluid-structure
interactions need to couple different types of equations together
and different parts of the FSI system behave differently so that the
resulting nonlinear algebraic system is hard to solve.
In this talk, we discuss a highly scalable, fully implicit multilevel
fluid-structure interaction solver based on isogeometric unstructured 3D
coarse spaces. The proposed
nonlinear solver is shown to be scalable with more than 10,000 processor cores.


\end{document}
